% !TEX root = ../thesis.tex

\chapter{Analýza požiadaviek}\label{ch:requirements-analysis}

Analýza požiadaviek sa zameriava na identifikáciu kľúčových oblastí navrhovaného informačného systému pre správu distribúcie vody, evidenciu revízií a servisných zásahov. Cieľom analýzy je previesť počiatočné textové požiadavky do formálnejšej podoby, ktorá bude použiteľná pri návrhu architektúry, databázového modelu a používateľských scenárov.

Výsledkom tejto fázy je systematické rozdelenie požiadaviek do štruktúrovaných kategórií, určenie ich priority a prepojenie na zdroje. Takto spracované požiadavky vytvárajú jednotný podklad pre implementáciu, testovanie a následné vyhodnotenie riešenia.

\section{Štruktúrovaná tabuľková reprezentácia požiadaviek}

Pri návrhu systému je dôležité, aby boli požiadavky definované jednoznačne, porovnateľne a overiteľne. Tabuľková reprezentácia umožňuje každú požiadavku presne identifikovať a priradiť jej vlastnosti, ktoré podporujú rozhodovanie o rozsahu implementácie.

V tejto kapitole je každá požiadavka opísaná prostredníctvom nasledujúcich atribútov:

\begin{itemize}
\item \textbf{ID}: jednoznačný identifikátor požiadavky.
\item \textbf{Popis požiadavky}: stručná formulácia očakávanej funkcionality.
\item \textbf{Klasifikácia}: zaradenie požiadavky do vecnej oblasti systému.
\item \textbf{Priorita}: dôležitosť požiadavky pre základnú prevádzku a implementačné poradie.
\item \textbf{Zdroj}: označenie pôvodu požiadavky (UR/SR) podľa jej charakteru.
\end{itemize}

Tabuľka \ref{tab:functional-requirements-structured} prezentuje funkčné požiadavky identifikované v predchádzajúcej kapitole a pripravuje podklad pre nadväzujúcu analýzu nefunkčných požiadaviek.

\subsection{Funkčné požiadavky}

\small
\begin{longtable}{|p{1.3cm}|p{5.8cm}|p{2.0cm}|c|p{2.1cm}|}
\caption{Štruktúrovaná tabuľková reprezentácia funkčných požiadaviek}\label{tab:functional-requirements-structured}\\
\hline
ID & Popis požiadavky & Klasifikácia & Priorita & Zdroj \\
\hline
\endfirsthead
\hline
ID & Popis požiadavky & Klasifikácia & Priorita & Zdroj \\
\hline
\endhead
FR-01 & Založenie karty zariadenia s identifikátorom, typom, umiestnením, stavom a dátumom uvedenia do prevádzky. & \shortstack[l]{Evidencia\\zariadení} & 5 & UR \\
\hline
FR-02 & Úprava údajov zariadenia s evidenciou histórie zmien. & \shortstack[l]{Evidencia\\zariadení} & 5 & UR \\
\hline
FR-03 & Vyhľadávanie a filtrovanie zariadení podľa typu, lokality, stavu a zodpovednej osoby. & \shortstack[l]{Evidencia\\zariadení} & 4 & UR \\
\hline
FR-04 & Plánovanie revízie zariadenia s termínom, rozsahom a pridelenou zodpovednou osobou. & Revízie & 5 & SR \\
\hline
FR-05 & Záznam výsledku revízie vrátane stavu, poznámky a príloh. & Revízie & 5 & SR \\
\hline
FR-06 & Zobrazenie histórie revízií zariadenia v časovom poradí. & Revízie & 4 & SR \\
\hline
FR-07 & Založenie servisného zásahu na základe poruchy, revízie alebo manuálne. & Servis & 5 & UR \\
\hline
FR-08 & Evidencia priebehu servisného zásahu vrátane vykonaných prác, použitého materiálu a zodpovednej osoby. & Servis & 4 & UR \\
\hline
FR-09 & Automatické odoslanie notifikácie pred plánovaným termínom revízie. & Notifikácie & 4 & SR \\
\hline
FR-10 & Automatické odoslanie notifikácie pri prekročení termínu revízie alebo servisného zásahu. & Notifikácie & 4 & SR \\
\hline
FR-11 & Generovanie reportu revízií za zvolené obdobie. & Reporty & 4 & UR \\
\hline
FR-12 & Export reportov minimálne do formátov PDF a CSV. & Reporty & 3 & SR \\
\hline
\end{longtable}

\noindent
Priorita je hodnotená na škále 1 až 5, kde 5 predstavuje kritickú požiadavku pre základnú prevádzku systému a 1 predstavuje požiadavku s najnižšou implementačnou prioritou.

\noindent
Označenia zdrojov: \textbf{UR} (User Requirements) a \textbf{SR} (System Requirements) vyjadrujú charakter požiadavky. \textbf{UR} predstavuje požiadavky odvodené z používateľských potrieb a prevádzkových procesov, kým \textbf{SR} predstavuje systémové požiadavky odvodené z legislatívnych a odborných podkladov, najmä \cite{EUDWD2020,SKVyhlaska912023,CarricoFerreira2021}.
